@ID: ON13D1P1U
@BIDANG: Struktur Aljabar

Suatu subgrup $H$ dari $G$ disebut subgrup normal jika untuk setiap $g\in G$ dan $h\in H$ berlaku $g^{-1}hg\in H$. Suatu subgrup $H$ dari $G$ dikatakan mempunyai sifat refleksif jika untuk setiap $a,b\in G$ dengan $ab\in H$ berlaku $ba\in H$.

Tunjukkan bahwa $H$ merupakan subgrup normal jika dan hanya jika $H$ bersifat refleksif.

@ID: ON13D1P2U
@BIDANG: Analisis Real

Misalkan $f:[0,1]\to[0,1]$ merupakan fungsi kontinu sedemikian sehingga

$$
\int_0^1f(x)\,dx=0.
$$

Tunjukkan bahwa untuk setiap $\alpha\in[0,1]$ berlaku

$$
\left|
\int_0^\alpha f(x)\,dx
\right|
\le
\frac18
\sup_{a\in[0,1]}|f'(a)|.
$$

@ID: ON13D1P3U
@BIDANG: Analisis Kompleks

Misalkan $u$ fungsi harmonik di $\mathbb R^2$ dan $v$ fungsi harmonik konjugatnya (yaitu fungsi sehingga $f(z)=u(z)+iv(z)$ analitik). Jika

$$
u^3-3uv^2\ge0
$$

di $\mathbb R^2$, tentukan semua fungsi $u$ dan $v$.

@ID: ON13D1P4U
@BIDANG: Aljabar Linear

Misalkan $A=[a_{ij}]$ matriks real berukuran $n\times n$ dengan

$$
a_{ij}=\min\{i,j\},
\qquad i,j=1,2,\ldots,n.
$$

Tentukan $A^{-1}$.

@ID: ON13D1P5U
@BIDANG: Kombinatorika

Sebuah $n$-board adalah sebuah persegi panjang berukuran $n\times1$.

Misalkan $f_n$ menyatakan banyaknya cara mengubin $n$-board dengan menggunakan ubin $1\times1$ dan ubin $2\times1$. Jadi $f_1=1$ dan $f_2=2$.

Untuk $n\ge0$, buktikan bahwa

$$
f_{2n+1}
=
\sum_{i\ge0}\sum_{j\ge0}
\binom{n-i}{j}
\binom{n-j}{i}.
$$

@ID: ON13D2P1U
@BIDANG: Kombinatorika

Untuk $0\le k\le n/2$, berikan bukti kombinatorial dari persamaan

$$
\sum_{m=k}^{n-k}
\binom{m}{k}
\binom{n-m}{k}
=
\binom{n+1}{2k+1},
$$

dengan $n$ bilangan bulat positif.

@ID: ON13D2P2U
@BIDANG: Struktur Aljabar

Misalkan $I$ suatu ideal di gelanggang $R$ dan $S\subseteq R$. Definisikan

$$
(I:S)=\{\,r\in R:rs\in I\text{ untuk setiap }s\in S\,\}.
$$

(a) Buktikan bahwa $(I:S)$ merupakan ideal kiri di $R$.

(b) Jika $S$ ideal di $R$, buktikan bahwa $(I:S)$ merupakan ideal di $R$.

(c) Jika $R$ daerah integral, $a,b\in R$ dengan $b\ne0$, dan $I=(ab)$, $J=(b)$, buktikan bahwa

$$
(I:J)=(a).
$$

@ID: ON13D2P3U
@BIDANG: Analisis Kompleks

Diketahui fungsi

$$
f(z)
=
\int_1^z
\left(
\frac1w+\frac{a}{w^3}
\right)
\cos w\,dw.
$$

Tentukan nilai $a$ agar fungsi tersebut bernilai tunggal di seluruh bidang kompleks.

@ID: ON13D2P4U
@BIDANG: Analisis Real

Misalkan $a_n\ge0$ untuk setiap $n\in\mathbb N$ dan $s_n=\sum_{k=1}^n a_k$. Jika $\sum_{n=1}^\infty a_n$ divergen, buktikan bahwa

$$
\sum_{n=1}^\infty
\frac{a_n}{(s_n)^\alpha}
$$

konvergen jika dan hanya jika $\alpha>1$.

@ID: ON13D2P5U
@BIDANG: Aljabar Linear

Misalkan $l$ dan $m$ dua garis di $\mathbb R^3$ yang bersilangan (yaitu tidak berpotongan dan tidak sejajar). Tentukan titik $P$ pada $l$ dan $Q$ pada $m$ yang meminimumkan jarak antara titik di $l$ dan titik di $m$.